(1) 取$\propto e^{i(k_xx+k_yy-\omega t)}$，得到$\begin{pmatrix}-i\omega&-f&igk_x\\f&-i\omega&igk_y\\iHk_x&iHk_y&-i\omega\end{pmatrix}(u,v,\eta)^T=0$。由于背景均匀且旋转固定，非零振幅要求行列式为零，从而得到三个线性本征频率。
(2) 对上述矩阵求行列式或先消去$(u,v)$，特征方程为$\omega[\omega^2-f^2-gH(k_x^2+k_y^2)]=0$。因此$\omega_\pm=\pm\sqrt{f^2+gHk^2}$是惯性重力波，第三支$\omega_0=0$是地转平衡模；后者存在不意味着全部扰动都能静止。
(3) 定义线性相对涡度$\zeta=v_x-u_y$。对动量方程分别求$x,y$导数并相减，得$\partial_t\zeta+f(u_x+v_y)=0$；联立连续性$\eta_t=-H(u_x+v_y)$得到$\partial_t(\zeta-f\eta/H)=0$。对零频模，速度与表面坡度满足$-fv=-g\eta_x$、$fu=-g\eta_y$，地转流沿等高线而非跨越坡度。
(4) 传播支群速度$\boldsymbol v_g=\nabla_{\boldsymbol k}\omega_\pm=\pm gH\boldsymbol k/\sqrt{f^2+gHk^2}$。在$k\to0$时频率趋$|f|$而群速度趋零，体现转动惯性振荡；在$k\to\infty$时$\omega_\pm\simeq\pm\sqrt{gH}\,k$，群速趋浅水重力波速$\sqrt{gH}$。零频地转模无此传播能流。
